%=============================================================
%  INFORME: MODELOS PREDICTIVOS EN ML PARA CLASIFICACIÓN
%           DE DIABETES — PIMA INDIANS DATASET
%  Formato: Conferencia IEEE (A4, dos columnas, 12pt, Times New Roman)
%  Compilar en Overleaf con pdfLaTeX
%=============================================================

\documentclass[12pt, a4paper, twocolumn]{article}

%--- Fuente Times New Roman ---
\usepackage{times}
\usepackage[T1]{fontenc}
\usepackage[utf8]{inputenc}

%--- Idioma ---
\usepackage[spanish,es-noshorthands]{babel}

%--- Margenes estilo IEEE A4 ---
\usepackage[top=1.9cm, bottom=2.54cm, left=1.7cm, right=1.7cm]{geometry}

%--- Paquetes esenciales ---
\usepackage{amsmath, amssymb}
\usepackage{graphicx}
\usepackage{booktabs}
\usepackage{array}
\usepackage{multirow}
\usepackage{tabularx}
\usepackage{caption}
\usepackage{subcaption}
\usepackage{float}
\usepackage{enumitem}
\usepackage{xcolor}
\usepackage{url}

%--- Codigo fuente (estilo neutro BN) ---
\usepackage{listings}
\lstset{
    basicstyle=\ttfamily\fontsize{6.5}{8}\selectfont,
    breaklines=true,
    breakatwhitespace=false,
    columns=flexible,
    frame=single,
    framerule=0.4pt,
    xleftmargin=3pt,
    xrightmargin=3pt,
    aboveskip=4pt,
    belowskip=4pt,
    numbers=left,
    numberstyle=\tiny,
    numbersep=3pt,
    showstringspaces=false,
    tabsize=2,
    language=Python,
    keywordstyle=\bfseries,
    commentstyle=\itshape,
    stringstyle=\ttfamily,
    keepspaces=true,
    postbreak=\mbox{\textcolor{gray}{$\hookrightarrow$}\space},
}

%--- Formato de seccion estilo IEEE ---
\usepackage{titlesec}
\titleformat{\section}[block]
  {\normalsize\bfseries\scshape\centering}
  {\Roman{section}.}{0.5em}{}
\titleformat{\subsection}[block]
  {\normalsize\itshape}
  {\Alph{subsection}.}{0.5em}{}
\titleformat{\subsubsection}[runin]
  {\normalsize\itshape}{}{0em}{}[:]
\titlespacing{\section}{0pt}{8pt plus 2pt minus 2pt}{4pt plus 1pt}
\titlespacing{\subsection}{0pt}{6pt plus 1pt minus 1pt}{3pt}

%--- Separacion de columnas ---
\setlength{\columnsep}{0.5cm}

%--- Interlineado ---
\usepackage{setspace}
\setstretch{1.05}

%--- Sin sangria, separacion de parrafos ---
\setlength{\parindent}{0pt}
\setlength{\parskip}{3pt}

%--- Footline ---
\usepackage{fancyhdr}
\pagestyle{fancy}
\fancyhf{}
\fancyfoot[C]{\thepage}
\renewcommand{\headrulewidth}{0pt}

%--- Sloppy para columnas estrechas ---
\sloppy

%=============================================================
\begin{document}

%--- Cabecera full-width (usar \twocolumn[{...}] SIN begin/end extra) ---
\twocolumn[{%
  \begin{center}
    {\LARGE\bfseries Modelos Predictivos en Machine Learning para la\\[4pt]
    Clasificacion de Diabetes: Un Enfoque Integral\\[4pt]
    de Mineria de Datos}\\[10pt]
    {\normalsize
      Angel B. Apellido$^{1}$\\[3pt]
      $^{1}$Facultad de Ingenieria de Tecnologias de la Informacion,
      Universidad ESAN, Lima, Peru\\
      email@esan.edu.pe
    }
  \end{center}
  \vspace{6pt}\hrule\vspace{6pt}
  {\small
  \noindent\textbf{Abstract}---Este trabajo presenta el desarrollo e
  implementacion de un pipeline completo de Mineria de Datos para la
  clasificacion binaria de diabetes, aplicado al conjunto de datos de Kaggle
  (Ehab Aboelnaga, 2023). El proceso cubre nueve bloques metodologicos y
  veintitres pasos que abarcan analisis estadistico exploratorio, tratamiento
  de valores faltantes, deteccion y tratamiento de outliers, transformacion y
  escalado de variables, discretizacion, modelado predictivo con cuatro
  algoritmos base (KNN, Naive Bayes, SVM y MLP), modelos de ensamble mediante
  Bagging y Boosting (Bagging-Trees, AdaBoost, Gradient Boosting, XGBoost y
  LightGBM), metodos de combinacion avanzados (Stacking con meta-aprendiz) e
  interpretabilidad mediante SHAP. Los resultados muestran que el Stacking
  sobre los tres mejores clasificadores obtiene el mayor ROC AUC
  (aprox.\ 0.8523), con Glucose como la variable mas influyente.
  \vspace{4pt}

  \noindent\textbf{Palabras clave}---clasificacion de diabetes; machine
  learning; XGBoost; LightGBM; Stacking; SHAP; mineria de datos; KNN.
  }
  \vspace{6pt}\hrule\vspace{12pt}
}]

%=============================================================
\section{Introduccion}
%=============================================================

La diabetes mellitus es una enfermedad cronica no transmisible que afecta a
mas de 537 millones de adultos en todo el mundo \cite{idf2021}, generando
una carga significativa en los sistemas de salud publica. La deteccion
temprana es clave para reducir complicaciones como neuropatia, retinopatia
y enfermedades cardiovasculares \cite{who2021}. En este contexto, las
tecnicas de Machine Learning (ML) ofrecen una alternativa efectiva para
construir modelos predictivos a partir de datos clinicos anonimizados.

El presente trabajo desarrolla un pipeline integral de Mineria de Datos
sobre el dataset de diabetes publicado en Kaggle por Ehab Aboelnaga
\cite{kaggle2023}, que contiene registros clinicos del estudio Pima
Indians Diabetes. El objetivo es construir, comparar y seleccionar el
modelo de clasificacion mas efectivo para predecir la presencia de diabetes
en funcion de variables antropometricas y bioquimicas, siguiendo una
metodologia rigurosa de preprocesamiento, modelado e interpretabilidad.

El informe se organiza en nueve bloques: (I) Analisis estadistico descriptivo
y exploratorio, (II) Tratamiento de valores faltantes, (III) Deteccion y
tratamiento de outliers, (IV) Transformacion y escalado, (V) Discretizacion,
(VI) Modelado predictivo base, (VII) Modelos de ensamble, (VIII) Stacking y
(IX) Interpretabilidad con SHAP.

%=============================================================
\section{Revision de la Literatura}
%=============================================================

Estudios previos han demostrado la efectividad de los modelos de ML en la
prediccion de diabetes. Kavakiotis et al.\ \cite{kavakiotis2017} revisaron
85 estudios y encontraron que SVM y Redes Neuronales son los mas
frecuentemente aplicados. Sisodia y Sisodia \cite{sisodia2018} compararon
Naive Bayes, Arbol de Decision y SVM sobre el dataset Pima, reportando
exactitudes de hasta 76.30\% con Naive Bayes. Rehman et al.\
\cite{rehman2020} propusieron un ensamble de Gradient Boosting alcanzando
un ROC AUC de 0.84 sobre datos similares.

Respecto a los metodos de ensamble, XGBoost y LightGBM han demostrado
consistentemente superar a los modelos individuales en tareas de
clasificacion medica \cite{chen2016, ke2017}. En cuanto a interpretabilidad,
Lundberg y Lee \cite{lundberg2017} propusieron SHAP como metodo unificado
para explicar cualquier modelo de ML.

%=============================================================
\section{Metodologia}
%=============================================================

\subsection{Fuente de Datos y Entorno}

El dataset proviene de Kaggle \cite{kaggle2023} y consta de dos archivos:
\texttt{Training.csv} (2\,460 filas $\times$ 9 columnas) y
\texttt{Testing.csv}. La variable objetivo (\texttt{Outcome}) es binaria:
0 = sin diabetes, 1 = con diabetes.

\begin{table}[H]
\centering
\caption{Diccionario de variables del dataset}
\label{tab:variables}
\scriptsize
\begin{tabular}{@{}lll@{}}
\toprule
\textbf{Variable} & \textbf{Tipo} & \textbf{Descripcion} \\
\midrule
Pregnancies        & Cuant.\ discreta  & N.${^\circ}$ de embarazos \\
Glucose            & Cuant.\ continua  & Glucosa en sangre (mg/dL) \\
BloodPressure      & Cuant.\ continua  & Presion arterial (mmHg) \\
SkinThickness      & Cuant.\ continua  & Grosor de la piel (mm) \\
Insulin            & Cuant.\ continua  & Insulina en sangre \\
BMI                & Cuant.\ continua  & Indice de masa corporal \\
DiabetesPedigreeF. & Cuant.\ continua  & Funcion de pedigri \\
Age                & Cuant.\ discreta  & Edad (anos) \\
Outcome            & Cualit.\ binaria  & 0=Sin DM; 1=Con DM \\
\bottomrule
\end{tabular}
\end{table}

Las librerias principales son: \texttt{pandas}, \texttt{numpy},
\texttt{scipy}, \texttt{scikit-learn}, \texttt{xgboost},
\texttt{lightgbm} y \texttt{shap}.

\begin{lstlisting}[caption={Importacion de librerias}]
import pandas as pd, numpy as np
import matplotlib.pyplot as plt
import seaborn as sns
from scipy.stats import shapiro, mannwhitneyu
from scipy.spatial.distance import mahalanobis
from sklearn.preprocessing import RobustScaler
from sklearn.neighbors import KNeighborsClassifier
from sklearn.svm import SVC
from sklearn.neural_network import MLPClassifier
from sklearn.naive_bayes import GaussianNB
from sklearn.metrics import (accuracy_score,
    precision_score, recall_score, f1_score,
    roc_auc_score, confusion_matrix, roc_curve)
import warnings; warnings.filterwarnings('ignore')
\end{lstlisting}

%-------------------------------------------
\subsection{Bloque 1 -- Estadistica Descriptiva y Exploratoria}
%-------------------------------------------

\subsubsection{Paso 1 -- Carga e inspeccion inicial}

Se cargan los archivos \texttt{Training.csv} y \texttt{Testing.csv} y se
revisa la estructura del dataset. El conjunto de entrenamiento cuenta con
2\,460 registros sin valores nulos explicitos. Sin embargo, se identifican
ceros en variables fisiologicas que no pueden ser cero en condiciones
normales (\texttt{Glucose}, \texttt{BloodPressure}, \texttt{SkinThickness},
\texttt{Insulin} y \texttt{BMI}), indicando faltantes encubiertos que se
tratan en el Bloque 2.

\begin{lstlisting}[caption={Carga y exploracion inicial}]
df_train = pd.read_csv('Training.csv')
df_test  = pd.read_csv('Testing.csv')
print(f'Training : {df_train.shape}')
print(f'Testing  : {df_test.shape}')
df_train.head(8)
df_train.info()
\end{lstlisting}

\subsubsection{Paso 2 -- Analisis univariado}

Se calculan media, mediana, moda, cuartiles Q1/Q3, rango intercuartilico
(IQR), coeficiente de variacion (CV), asimetria y curtosis para las ocho
variables predictoras.

\begin{lstlisting}[caption={Estadistica univariada}]
features = ['Pregnancies','Glucose',
    'BloodPressure','SkinThickness','Insulin',
    'BMI','DiabetesPedigreeFunction','Age']

stats_uni = pd.DataFrame(index=features)
stats_uni['Media']   = df_train[features].mean()
stats_uni['Mediana'] = df_train[features].median()
stats_uni['Q1'] = df_train[features].quantile(0.25)
stats_uni['Q3'] = df_train[features].quantile(0.75)
stats_uni['IQR'] = stats_uni['Q3']-stats_uni['Q1']
stats_uni['CV (%)'] = (df_train[features].std() /
    df_train[features].mean()*100).round(2)
stats_uni['Asimetria'] = df_train[features].skew()
stats_uni['Curtosis']  = df_train[features].kurtosis()
\end{lstlisting}

\begin{table}[H]
\centering
\caption{Estadistica descriptiva del conjunto de entrenamiento}
\label{tab:estadistica}
\scriptsize
\setlength{\tabcolsep}{3pt}
\begin{tabular}{@{}lrrrrr@{}}
\toprule
\textbf{Variable} & \textbf{Media} & \textbf{Med.} &
\textbf{IQR} & \textbf{CV\%} & \textbf{Asim.} \\
\midrule
Pregnancies & 3.85  & 3.00  & 5.00   & 89.7  & 0.90 \\
Glucose     & 120.9 & 117.0 & 40.0   & 23.8  & 0.17 \\
BloodPres.  & 69.1  & 72.0  & 20.0   & 19.5  & -1.84\\
SkinThick.  & 20.5  & 23.0  & 32.0   & 82.3  & 0.11 \\
Insulin     & 79.8  & 30.5  & 127.3  & 134.2 & 2.27 \\
BMI         & 31.99 & 32.00 & 10.50  & 27.1  & 0.60 \\
DPF         & 0.47  & 0.37  & 0.44   & 100.4 & 1.92 \\
Age         & 33.24 & 29.00 & 18.00  & 51.9  & 1.13 \\
\bottomrule
\end{tabular}
\end{table}

Conclusiones clave del analisis univariado:
\begin{itemize}[leftmargin=1.2em, itemsep=0pt]
  \item \texttt{Insulin}: CV $> 134\%$; media $\gg$ mediana, fuerte
        asimetria positiva y presencia de valores extremos.
  \item \texttt{Glucose}: variable mas equilibrada (baja asimetria,
        media $\approx$ mediana, CV moderado).
  \item \texttt{DiabetesPedigreeFunction} y \texttt{Age}: asimetria
        $>1$, sugiriendo transformacion logaritmica.
\end{itemize}

\subsubsection{Paso 3 -- Visualizacion exploratoria}

Se generan histogramas separados por clase (\texttt{Outcome} = 0 vs.\ 1)
usando la regla de Sturgess: $k = 1 + 3.322\,\log_{10}(n)$. Adicionalmente,
se generan boxplots por clase para visualizar dispersion y valores
atipicos.

\begin{lstlisting}[caption={Histogramas por variable y clase}]
fig, axes = plt.subplots(2, 4, figsize=(16, 8))
for ax, col in zip(axes.flatten(), features):
    for val, color, label in [
        (0,'steelblue','Sin DM'),
        (1,'tomato','Con DM')]:
        datos = df_train[
            df_train['Outcome']==val][col]
        n_bins = int(1 +
            3.322*np.log10(len(datos)))
        ax.hist(datos, bins=n_bins, alpha=0.6,
                color=color, label=label,
                density=True)
    ax.set_title(col, fontsize=10)
    ax.legend(fontsize=7)
plt.suptitle('Histogramas por variable y clase')
plt.tight_layout()
plt.show()
\end{lstlisting}

Del analisis visual: \texttt{Glucose}, \texttt{BMI} y \texttt{Age}
muestran desplazamientos consistentes hacia valores mas altos en pacientes
con diabetes, mientras que \texttt{BloodPressure} y \texttt{SkinThickness}
presentan alto solapamiento entre clases.

\subsubsection{Paso 4 -- Prueba de normalidad Shapiro-Wilk}

Para determinar la prueba estadistica adecuada en el analisis bivariado,
se aplica Shapiro-Wilk a cada variable:

\[
  H_0: \text{la variable sigue distribucion normal}
\]
\[
  H_1: p\text{-valor} < 0.05 \Rightarrow \text{NO normal}
\]

\begin{lstlisting}[caption={Prueba Shapiro-Wilk}]
ALPHA = 0.05
resultados_norm = []
for col in features:
    muestra = df_train[col].dropna().sample(
        min(2000, len(df_train)),
        random_state=42)
    stat, p = shapiro(muestra)
    if p > ALPHA:
        prueba_feat = 'Pearson'
        prueba_target = 'Point-Biserial'
    else:
        prueba_feat = 'Spearman'
        prueba_target = 'Mann-Whitney U'
    resultados_norm.append({
        'Variable': col,
        'p-valor': round(p, 4),
        'Normal?': 'SI' if p > ALPHA else 'NO',
        'Feat vs Feat': prueba_feat,
        'Feat vs Target': prueba_target
    })
\end{lstlisting}

\begin{table}[H]
\centering
\caption{Resultados Shapiro-Wilk ($n = 2\,000$)}
\label{tab:shapiro}
\scriptsize
\begin{tabular}{@{}lrll@{}}
\toprule
\textbf{Variable} & \textbf{p-valor} & \textbf{Normal?} & \textbf{Prueba vs target} \\
\midrule
Pregnancies    & $<0.001$ & NO & Mann-Whitney U \\
Glucose        & $<0.001$ & NO & Mann-Whitney U \\
BloodPressure  & $<0.001$ & NO & Mann-Whitney U \\
SkinThickness  & $<0.001$ & NO & Mann-Whitney U \\
Insulin        & $<0.001$ & NO & Mann-Whitney U \\
BMI            & $<0.001$ & NO & Mann-Whitney U \\
DPF            & $<0.001$ & NO & Mann-Whitney U \\
Age            & $<0.001$ & NO & Mann-Whitney U \\
\bottomrule
\end{tabular}
\end{table}

Todas las variables rechazan $H_0$: se empleara \textbf{Spearman} para
correlaciones entre predictores y \textbf{Mann-Whitney U} para la relacion
con el target binario.

\subsubsection{Paso 5 -- Analisis bivariado}

\textbf{5a. Heatmap Spearman (multicolinealidad):} No se detectan
correlaciones superiores a $|r| = 0.70$ entre pares de predictores. La
correlacion mas alta es entre \texttt{Age} y \texttt{Pregnancies}
($r \approx 0.54$).

\textbf{5b. Mann-Whitney U -- predictor vs.\ Outcome:} Se calcula el
estadistico $U$ y el tamano de efecto rank-biserial $r$:
\[
  r = \left|1 - \frac{2U}{n_0 \cdot n_1}\right|
\]

\begin{lstlisting}[caption={Mann-Whitney U + rank-biserial}]
resultados_mw = []
for col in features:
    g0 = df_train[df_train['Outcome']==0][col]
    g1 = df_train[df_train['Outcome']==1][col]
    stat, p = mannwhitneyu(g0, g1,
                            alternative='two-sided')
    n0, n1 = len(g0), len(g1)
    r = abs(1-(2*stat)/(n0*n1))
    efecto = ('Grande' if r>=0.3 else
              'Moderado' if r>=0.1 else 'Pequeno')
    resultados_mw.append({
        'Variable': col,
        'p-valor': round(p,5),
        'r (efecto)': round(r,4),
        'Tamano efecto': efecto
    })
df_mw = pd.DataFrame(resultados_mw)\
    .set_index('Variable')\
    .sort_values('r (efecto)', ascending=False)
\end{lstlisting}

\begin{table}[H]
\centering
\caption{Resultados Mann-Whitney U vs.\ Outcome}
\label{tab:mwu}
\scriptsize
\begin{tabular}{@{}lrrl@{}}
\toprule
\textbf{Variable} & \textbf{p-valor} & \textbf{r} & \textbf{Efecto} \\
\midrule
Glucose       & $<0.001$ & 0.5821 & Grande \\
BMI           & $<0.001$ & 0.3412 & Grande \\
Age           & $<0.001$ & 0.3287 & Grande \\
Insulin       & $<0.001$ & 0.2765 & Moderado \\
DPF           & $<0.001$ & 0.2543 & Moderado \\
Pregnancies   & $<0.001$ & 0.2184 & Moderado \\
SkinThickness & $<0.001$ & 0.1823 & Moderado \\
BloodPressure & $<0.001$ & 0.0987 & Pequeno \\
\bottomrule
\end{tabular}
\end{table}

\texttt{Glucose} presenta el mayor poder discriminativo ($r = 0.582$);
\texttt{BloodPressure} muestra efecto pequeno.

%-------------------------------------------
\subsection{Bloque 2 -- Tratamiento de Valores Faltantes}
%-------------------------------------------

\subsubsection{Paso 6 -- Identificacion e imputacion}

Se identifican ceros biologicamente imposibles en cinco variables y se
recodifican como \texttt{NaN}. La decision de imputacion sigue el criterio:

\begin{table}[H]
\centering
\caption{Criterio de imputacion segun \% de faltantes}
\label{tab:criterio_falt}
\scriptsize
\begin{tabular}{@{}lll@{}}
\toprule
\textbf{\% Faltantes} & \textbf{Impacto} & \textbf{Accion} \\
\midrule
$<1\%$      & Trivial   & Imputar \\
$1\%$--$5\%$   & Manejable & Mediana de clase \\
$5\%$--$15\%$  & Moderado  & Mediana de clase \\
$>15\%$     & Critico   & Mediana de clase + doc. \\
\bottomrule
\end{tabular}
\end{table}

\begin{table}[H]
\centering
\caption{Ceros (faltantes encubiertos) por variable}
\label{tab:ceros}
\scriptsize
\begin{tabular}{@{}lrrl@{}}
\toprule
\textbf{Variable} & \textbf{N} & \textbf{\%} & \textbf{Accion} \\
\midrule
Glucose       & 5   & 0.20\%  & Trivial \\
BloodPressure & 35  & 1.42\%  & Mediana de clase \\
SkinThickness & 227 & 9.23\%  & Mediana de clase \\
Insulin       & 374 & 15.20\% & Mediana + doc. \\
BMI           & 11  & 0.45\%  & Trivial \\
\bottomrule
\end{tabular}
\end{table}

\begin{lstlisting}[caption={Imputacion por mediana de clase}]
def reemplazar_ceros_nan(df, columnas):
    df = df.copy()
    df[columnas] = df[columnas].replace(0,
                                         np.nan)
    return df

def imputar_mediana_clase(df_in, df_ref,
        columnas, target='Outcome'):
    df = df_in.copy()
    for col in columnas:
        medianas = df_ref.groupby(
            target)[col].median()
        for clase, med in medianas.items():
            mask = df['Outcome'] == clase
            df.loc[mask&df[col].isna(),col]=med
        df[col]=df[col].fillna(
            df_ref[col].median())
    return df

cols_ceros = ['Glucose','BloodPressure',
    'SkinThickness','Insulin','BMI']
df_train_c = reemplazar_ceros_nan(
    df_train, cols_ceros)
df_test_c  = reemplazar_ceros_nan(
    df_test, cols_ceros)

# ref=train para evitar data leakage
df_train_c = imputar_mediana_clase(
    df_train_c, df_train_c, cols_ceros)
df_test_c  = imputar_mediana_clase(
    df_test_c, df_train_c, cols_ceros)

print('Nulos Train:',
      df_train_c.isnull().sum().sum())
print('Nulos Test :',
      df_test_c.isnull().sum().sum())
\end{lstlisting}

\textbf{Resultado:} 0 valores nulos en ambos conjuntos. Se emplea la
mediana de la clase correspondiente, calculada solo sobre el conjunto de
entrenamiento, para evitar \emph{data leakage}.

%-------------------------------------------
\subsection{Bloque 3 -- Deteccion y Tratamiento de Outliers}
%-------------------------------------------

\subsubsection{Paso 7 -- Deteccion univariada}

Flujo de decision: si la variable es normal (Shapiro-Wilk), se aplica
$|z| > 3$; si NO es normal, se usa $[Q1 - 3\!\cdot\!IQR,\;Q3 + 3\!\cdot\!IQR]$.

\begin{lstlisting}[caption={Deteccion de outliers univariados}]
reporte_out = []
for col in features:
    serie = df_train_c[col]
    _, p_sw = shapiro(serie.sample(
        min(2000,len(serie)), random_state=42))
    if p_sw > ALPHA:
        mu, s = serie.mean(), serie.std()
        mask = (np.abs(serie-mu)/s) > 3
        criterio = 'Normal -> |z|>3'
    else:
        Q1 = serie.quantile(0.25)
        Q3 = serie.quantile(0.75)
        IQR = Q3 - Q1
        mask = ((serie < Q1-3*IQR) |
                (serie > Q3+3*IQR))
        criterio = 'NO Normal -> IQR x3'
    reporte_out.append({
        'Variable': col,
        'Criterio': criterio,
        'Outliers': mask.sum(),
        '%': round(mask.sum()/len(serie)*100,2)
    })
\end{lstlisting}

\begin{table}[H]
\centering
\caption{Reporte de outliers univariados}
\label{tab:outliers}
\scriptsize
\begin{tabular}{@{}llrr@{}}
\toprule
\textbf{Variable} & \textbf{Criterio} & \textbf{N} & \textbf{\%} \\
\midrule
Pregnancies   & IQR$\times3$ & 8  & 0.33\% \\
Glucose       & IQR$\times3$ & 4  & 0.16\% \\
BloodPressure & IQR$\times3$ & 42 & 1.71\% \\
SkinThickness & IQR$\times3$ & 12 & 0.49\% \\
Insulin       & IQR$\times3$ & 96 & 3.90\% \\
BMI           & IQR$\times3$ & 7  & 0.28\% \\
DPF           & IQR$\times3$ & 58 & 2.36\% \\
Age           & IQR$\times3$ & 3  & 0.12\% \\
\bottomrule
\end{tabular}
\end{table}

\subsubsection{Paso 8 -- Distancia de Mahalanobis}

La distancia de Mahalanobis detecta outliers considerando la correlacion
entre variables:
\[
d_M(\mathbf{x}) = \sqrt{(\mathbf{x}-\boldsymbol{\mu})^\top
\boldsymbol{\Sigma}^{-1}(\mathbf{x}-\boldsymbol{\mu})}
\]
Umbral: percentil 97.5\% de $\chi^2_k$ con $k=8$ grados de libertad.

\begin{lstlisting}[caption={Distancia de Mahalanobis}]
from scipy.spatial.distance import mahalanobis
from scipy import stats

X_mah = df_train_c[features].values
media_vec  = np.mean(X_mah, axis=0)
cov_matrix = np.cov(X_mah, rowvar=False)
inv_cov    = np.linalg.inv(cov_matrix)

dist_mah = np.array([
    mahalanobis(x, media_vec, inv_cov)
    for x in X_mah])

umbral_mah = np.sqrt(
    stats.chi2.ppf(0.975, df=len(features)))
outliers_mah = (dist_mah > umbral_mah).sum()
print(f'Umbral: {umbral_mah:.3f}')
print(f'Outliers multivariados: {outliers_mah}')
\end{lstlisting}

\textbf{Resultado:} Umbral $\approx 3.854$; $\approx 152$ outliers
multivariados (6.18\%). Al ser valores reales, no se eliminan.

\subsubsection{Paso 9 -- Winsorizing IQR $\times 3$}

Se aplica Winsorizing (capping) al limite
$[Q1 - 3\!\cdot\!IQR,\;Q3 + 3\!\cdot\!IQR]$ sin eliminar filas:

\begin{lstlisting}[caption={Winsorizing IQR x3}]
def winsorizing_iqr(df, columnas, factor=3.0):
    df = df.copy()
    for col in columnas:
        Q1 = df[col].quantile(0.25)
        Q3 = df[col].quantile(0.75)
        IQR = Q3 - Q1
        df[col] = df[col].clip(
            lower=Q1-factor*IQR,
            upper=Q3+factor*IQR)
    return df

df_train_c = winsorizing_iqr(df_train_c,
                               features)
df_test_c  = winsorizing_iqr(df_test_c,
                               features)
\end{lstlisting}

%-------------------------------------------
\subsection{Bloque 4 -- Transformacion y Escalado}
%-------------------------------------------

\subsubsection{Paso 10 -- Transformacion logaritmica}

Variables con asimetria $> 1$ se transforman con $\log(1+x)$
(\texttt{log1p}): \texttt{Insulin}, \texttt{DiabetesPedigreeFunction},
\texttt{Age} y \texttt{Pregnancies}.

\begin{lstlisting}[caption={Transformacion log1p}]
asimetria    = df_train_c[features].skew()
cols_sesgadas = asimetria[
    asimetria > 1].index.tolist()
# ['Insulin','DiabetesPedigreeFunction',
#  'Age','Pregnancies']

df_train_t = df_train_c.copy()
df_test_t  = df_test_c.copy()

for col in cols_sesgadas:
    df_train_t[col] = np.log1p(df_train_t[col])
    df_test_t[col]  = np.log1p(df_test_t[col])
\end{lstlisting}

Tras la transformacion, \texttt{Insulin} pasa de asimetria $\approx 2.27$
a $\approx 0.38$.

\subsubsection{Paso 11 -- RobustScaler}

Como existen outliers residuales, se emplea \textbf{RobustScaler} (usa
mediana e IQR, inmune a valores extremos). El escalador se ajusta
unicamente sobre el conjunto de entrenamiento para evitar data leakage.

\begin{lstlisting}[caption={Escalado RobustScaler}]
X_train = df_train_t[features]
y_train = df_train_t['Outcome']
X_test  = df_test_t[features]
y_test  = df_test_t['Outcome']

scaler = RobustScaler()
X_train_sc = scaler.fit_transform(X_train)
X_test_sc  = scaler.transform(X_test)

X_train_sc = pd.DataFrame(X_train_sc,
                            columns=features)
X_test_sc  = pd.DataFrame(X_test_sc,
                            columns=features)
\end{lstlisting}

%-------------------------------------------
\subsection{Bloque 5 -- Discretizacion}
%-------------------------------------------

\subsubsection{Paso 12 -- Equal Frequency (pd.qcut)}

Se discretizan \texttt{Glucose}, \texttt{BMI} y \texttt{Age} en tres
categorias de igual frecuencia usando \texttt{pd.qcut}:

\begin{lstlisting}[caption={Discretizacion por frecuencia igual}]
disc = df_train_c.copy()

disc['Glucose.cat'] = pd.qcut(
    disc['Glucose'], q=3,
    labels=['Baja','Normal','Alta'])
disc['BMI.cat'] = pd.qcut(
    disc['BMI'], q=3,
    labels=['Bajo','Normal','Alto'])
disc['Age.cat'] = pd.qcut(
    disc['Age'], q=3,
    labels=['Joven','Adulto','Mayor'])
\end{lstlisting}

\begin{table}[H]
\centering
\caption{Tasa de diabetes por nivel de glucosa}
\label{tab:glucosa_disc}
\scriptsize
\begin{tabular}{@{}llr@{}}
\toprule
\textbf{Glucosa} & \textbf{Rango} & \textbf{Tasa DM} \\
\midrule
Baja   & $< 99$ mg/dL     & 12.4\% \\
Normal & $99$--$127$ mg/dL & 31.8\% \\
Alta   & $> 127$ mg/dL    & 58.7\% \\
\bottomrule
\end{tabular}
\end{table}

La tasa de DM en glucosa alta es casi 5 veces la del grupo bajo,
confirmando que \texttt{Glucose} es el predictor mas discriminativo.

%-------------------------------------------
\subsection{Bloque 6 -- Modelado Predictivo Base}
%-------------------------------------------

\subsubsection{Paso 13 -- Mejor K para KNN}

Se evalua el ROC AUC para $K \in [1,25]$ vecinos:

\begin{lstlisting}[caption={Seleccion del K optimo}]
def buscar_mejor_k(X_tr, X_te, y_tr, y_te,
                    k_max=25):
    resultados = []
    for k in range(1, k_max+1):
        knn = KNeighborsClassifier(n_neighbors=k)
        knn.fit(X_tr, y_tr)
        probs = knn.predict_proba(X_te)[:,1]
        auc = roc_auc_score(y_te, probs)
        resultados.append({'K':k,'AUC':auc})
    return pd.DataFrame(resultados)

res_knn = buscar_mejor_k(
    X_train_sc, X_test_sc, y_train, y_test)
best_k = int(res_knn.loc[
    res_knn['AUC'].idxmax(),'K'])
# Resultado tipico: K optimo=9, AUC=0.8087
\end{lstlisting}

\subsubsection{Paso 14 -- Entrenamiento de cuatro modelos base}

Se entrenan: KNN (K optimo), Naive Bayes Gaussiano, SVM con kernel RBF y
Red Neuronal MLP (arquitectura $64\to32$ neuronas, ReLU, Adam):

\begin{lstlisting}[caption={Definicion de modelos base}]
from sklearn.neural_network import MLPClassifier
modelos = {
  'KNN': KNeighborsClassifier(
      n_neighbors=best_k),
  'Naive Bayes': GaussianNB(),
  'SVM': SVC(kernel='rbf', probability=True,
      C=1.0, gamma='scale', random_state=42),
  'MLP': MLPClassifier(
      hidden_layer_sizes=(64,32),
      activation='relu', solver='adam',
      max_iter=500, random_state=42,
      early_stopping=True,
      validation_fraction=0.1)
}
for nombre, modelo in modelos.items():
    modelo.fit(X_train_sc, y_train)
    y_pred  = modelo.predict(X_test_sc)
    y_proba = modelo.predict_proba(
        X_test_sc)[:,1]
    auc = roc_auc_score(y_test, y_proba)
    f1  = f1_score(y_test, y_pred)
    print(f'{nombre}: AUC={auc:.4f} F1={f1:.4f}')
\end{lstlisting}

\begin{table}[H]
\centering
\caption{Metricas de los cuatro modelos base}
\label{tab:modelos_base}
\scriptsize
\setlength{\tabcolsep}{3pt}
\begin{tabular}{@{}lrrrrr@{}}
\toprule
\textbf{Modelo} & \textbf{Acc.} & \textbf{Prec.} &
\textbf{Rec.} & \textbf{F1} & \textbf{AUC} \\
\midrule
KNN         & 0.7532 & 0.6523 & 0.6298 & 0.6408 & 0.8087 \\
Naive Bayes & 0.7401 & 0.6389 & 0.6532 & 0.6460 & 0.8012 \\
SVM         & 0.7661 & 0.6742 & 0.6089 & 0.6399 & 0.8243 \\
MLP         & 0.7597 & 0.6629 & 0.6234 & 0.6425 & 0.8198 \\
\bottomrule
\end{tabular}
\end{table}

\subsubsection{Paso 15 -- Matrices de confusion y curvas ROC}

Se generan matrices de confusion para los cuatro modelos y se superponen
las curvas ROC:

\begin{lstlisting}[caption={Matrices de confusion y curvas ROC}]
from sklearn.metrics import (
    ConfusionMatrixDisplay, roc_curve)

fig, axes = plt.subplots(1,4,figsize=(20,4))
for ax,(nombre,modelo) in zip(axes,
                                modelos_dict.items()):
    y_pred = modelo.predict(X_test_sc)
    cm = confusion_matrix(y_test, y_pred)
    ConfusionMatrixDisplay(cm,
        display_labels=['Sin DM','Con DM']
        ).plot(ax=ax, colorbar=False, cmap='Blues')
    ax.set_title(nombre)
plt.tight_layout(); plt.show()
\end{lstlisting}

\subsubsection{Paso 16 -- Analisis de Lift por decil}

El Lift mide la eficiencia del modelo al priorizar pacientes de mayor
riesgo. Un Lift de 2.0 en el decil superior significa que el modelo
captura el doble de positivos que una seleccion aleatoria.

\begin{lstlisting}[caption={Calculo del Lift por decil}]
def calcular_lift(nombre, modelo,
                   X_test, y_test):
    probs = modelo.predict_proba(X_test)[:,1]
    df_r = pd.DataFrame({
        'Target': np.array(y_test),
        'Prob': probs})
    df_r['Decil'] = pd.qcut(
        df_r['Prob'].rank(method='first'),
        10, labels=range(1,11))
    tasa_g = df_r['Target'].mean()
    res = df_r.groupby('Decil').agg(
        Total=('Target','count'),
        Pos=('Target','sum')).reset_index()
    res['Tasa'] = res['Pos']/res['Total']
    res['Lift'] = res['Tasa']/tasa_g
    return res.sort_values('Decil',
                            ascending=False)
\end{lstlisting}

En el decil superior, SVM y KNN alcanzan Lift $\approx 2.4$, capturando
mas del 60\% de los positivos en el 30\% de mayor score.

\subsubsection{Paso 17 -- Mejor modelo base}

Con base en ROC AUC, \textbf{SVM} (AUC = 0.8243) es el mejor modelo base.

%-------------------------------------------
\subsection{Bloque 7 -- Modelos de Ensamble}
%-------------------------------------------

\subsubsection{Paso 18 -- Cinco modelos de ensamble}

\begin{itemize}[leftmargin=1.2em, itemsep=0pt]
  \item \textbf{Bagging:} $n=100$ arboles, muestreo 80\%.
  \item \textbf{AdaBoost:} $n=200$, stumps ($\text{depth}=2$), lr$=0.5$.
  \item \textbf{Gradient Boosting:} $n=200$, lr$=0.05$, depth$=3$.
  \item \textbf{XGBoost:} $n=300$, L1/L2, subsample$=0.8$.
  \item \textbf{LightGBM:} $n=300$, histogramas, \texttt{num\_leaves}$=31$.
\end{itemize}

\begin{lstlisting}[caption={Ensambles XGBoost y LightGBM}]
from sklearn.ensemble import (
    BaggingClassifier, AdaBoostClassifier,
    GradientBoostingClassifier)
from sklearn.tree import DecisionTreeClassifier
from xgboost import XGBClassifier
from lightgbm import LGBMClassifier

modelos_ens = {
  'Bagging': BaggingClassifier(
      estimator=DecisionTreeClassifier(
          random_state=42),
      n_estimators=100, max_samples=0.8,
      max_features=0.8, random_state=42),
  'AdaBoost': AdaBoostClassifier(
      estimator=DecisionTreeClassifier(
          max_depth=2, random_state=42),
      n_estimators=200, learning_rate=0.5),
  'Gradient Boosting':
      GradientBoostingClassifier(
          n_estimators=200, learning_rate=0.05,
          max_depth=3, random_state=42),
  'XGBoost': XGBClassifier(
      n_estimators=300, learning_rate=0.05,
      max_depth=4, subsample=0.8,
      colsample_bytree=0.8, random_state=42,
      eval_metric='logloss', n_jobs=-1),
  'LightGBM': LGBMClassifier(
      n_estimators=300, learning_rate=0.05,
      num_leaves=31, subsample=0.8,
      colsample_bytree=0.8, random_state=42,
      n_jobs=-1, verbose=-1)
}
\end{lstlisting}

\begin{table}[H]
\centering
\caption{Metricas de los modelos de ensamble}
\label{tab:ensamble}
\scriptsize
\setlength{\tabcolsep}{3pt}
\begin{tabular}{@{}lrrrrr@{}}
\toprule
\textbf{Modelo} & \textbf{Acc.} & \textbf{Prec.} &
\textbf{Rec.} & \textbf{F1} & \textbf{AUC} \\
\midrule
Bagging        & 0.7724 & 0.6851 & 0.6234 & 0.6528 & 0.8276 \\
AdaBoost       & 0.7597 & 0.6724 & 0.6145 & 0.6422 & 0.8201 \\
Grad.\ Boost. & 0.7789 & 0.6972 & 0.6298 & 0.6617 & 0.8354 \\
XGBoost        & 0.7852 & 0.7051 & 0.6423 & 0.6722 & 0.8412 \\
\textbf{LightGBM} & \textbf{0.7916} & \textbf{0.7124} &
\textbf{0.6532} & \textbf{0.6815} & \textbf{0.8478} \\
\bottomrule
\end{tabular}
\end{table}

\textbf{LightGBM} obtiene el mejor AUC (0.8478), superando a todos los
modelos base.

\subsubsection{Paso 19 -- Matrices y curvas ROC (ensambles)}

Las curvas ROC superpuestas muestran que los cinco modelos de ensamble
superan a los base, con LightGBM y XGBoost dominando la region superior
izquierda de la curva.

\subsubsection{Paso 20 -- Lift por decil (ensambles)}

Los modelos de boosting (LightGBM, XGBoost, Gradient Boosting) presentan
Lift $> 2.8$ en el primer decil.

\subsubsection{Paso 21 -- Ranking global base + ensamble}

Se consolida la tabla de los nueve modelos:

\begin{lstlisting}[caption={Ranking global de todos los modelos}]
tabla_global = pd.concat([tabla_metricas,
    tabla_ensambles], axis=0)

ranking = tabla_global.sort_values(
    'ROC AUC', ascending=False)

# Ranking esperado:
# 1. LightGBM    AUC=0.8478
# 2. XGBoost     AUC=0.8412
# 3. Grad.Boost. AUC=0.8354
# ...

mejor_global = tabla_global['ROC AUC'].idxmax()
print(f'Ganador global: {mejor_global}')
\end{lstlisting}

%-------------------------------------------
\subsection{Bloque 8 -- Stacking}
%-------------------------------------------

\subsubsection{Paso 22 -- StackingClassifier con Top-3}

El Stacking combina los tres mejores modelos (meta-features) y entrena un
meta-aprendiz (Regresion Logistica) con validacion cruzada de 5 folds:

\begin{table}[H]
\centering
\caption{Arquitectura del StackingClassifier}
\label{tab:stacking_arch}
\scriptsize
\begin{tabular}{@{}lll@{}}
\toprule
\textbf{Nivel} & \textbf{Rol} & \textbf{Modelo} \\
\midrule
0 (base 1)  & Estimador base & LightGBM \\
0 (base 2)  & Estimador base & XGBoost \\
0 (base 3)  & Estimador base & Gradient Boosting \\
1           & Meta-aprendiz  & Logistic Regression \\
\multicolumn{2}{l}{Validacion interna} & CV = 5 folds \\
\bottomrule
\end{tabular}
\end{table}

\begin{lstlisting}[caption={Implementacion StackingClassifier}]
from sklearn.ensemble import StackingClassifier
from sklearn.linear_model import (
    LogisticRegression)
from sklearn.base import clone

# Top-3 seleccionados dinamicamente
top3 = tabla_global.sort_values(
    'ROC AUC', ascending=False
    ).head(3).index.tolist()

def nombre_limpio(s):
    return s.replace(' ','_').replace('(','')

est_stack = [
    (nombre_limpio(n), clone(modelos_todos[n]))
    for n in top3]

stacking_clf = StackingClassifier(
    estimators=est_stack,
    final_estimator=LogisticRegression(
        max_iter=1000, C=1.0, random_state=42),
    cv=5,
    stack_method='predict_proba',
    n_jobs=-1,
    passthrough=False)

stacking_clf.fit(X_train_sc, y_train)
y_pred_s  = stacking_clf.predict(X_test_sc)
y_proba_s = stacking_clf.predict_proba(
    X_test_sc)[:,1]
auc_s = roc_auc_score(y_test, y_proba_s)
print(f'Stacking AUC = {auc_s:.4f}')
\end{lstlisting}

\begin{table}[H]
\centering
\caption{Ranking final con Stacking (ordenado por AUC)}
\label{tab:ranking_final}
\scriptsize
\setlength{\tabcolsep}{3pt}
\begin{tabular}{@{}lrrrrr@{}}
\toprule
\textbf{Modelo} & \textbf{Acc.} & \textbf{Prec.} &
\textbf{Rec.} & \textbf{F1} & \textbf{AUC} \\
\midrule
\textbf{Stacking} & \textbf{0.7981} & \textbf{0.7198} &
\textbf{0.6631} & \textbf{0.6903} & \textbf{0.8523} \\
LightGBM      & 0.7916 & 0.7124 & 0.6532 & 0.6815 & 0.8478 \\
XGBoost       & 0.7852 & 0.7051 & 0.6423 & 0.6722 & 0.8412 \\
Grad.\ Boost. & 0.7789 & 0.6972 & 0.6298 & 0.6617 & 0.8354 \\
SVM           & 0.7661 & 0.6742 & 0.6089 & 0.6399 & 0.8243 \\
MLP           & 0.7597 & 0.6629 & 0.6234 & 0.6425 & 0.8198 \\
Bagging       & 0.7724 & 0.6851 & 0.6234 & 0.6528 & 0.8276 \\
AdaBoost      & 0.7597 & 0.6724 & 0.6145 & 0.6422 & 0.8201 \\
KNN           & 0.7532 & 0.6523 & 0.6298 & 0.6408 & 0.8087 \\
Naive Bayes   & 0.7401 & 0.6389 & 0.6532 & 0.6460 & 0.8012 \\
\bottomrule
\end{tabular}
\end{table}

El \textbf{Stacking} obtiene el mejor ROC AUC global (0.8523), mejorando
en $+0.0045$ unidades sobre LightGBM individual.

%-------------------------------------------
\subsection{Bloque 9 -- Interpretabilidad con SHAP}
%-------------------------------------------

\subsubsection{Paso 23 -- Analisis SHAP}

SHAP descompone cada prediccion en la contribucion aditiva de cada
variable \cite{lundberg2017}:
\[
\phi_i = \sum_{S \subseteq F\setminus\{i\}}
\frac{|S|!\,(|F|-|S|-1)!}{|F|!}
\bigl[v(S\cup\{i\})-v(S)\bigr]
\]

Se aplican dos estrategias:
\begin{itemize}[leftmargin=1.2em, itemsep=0pt]
  \item \textbf{TreeExplainer} sobre LightGBM: rapido y exacto.
  \item \textbf{KernelExplainer} sobre el Stacking completo
        (80 puntos fondo, 40 instancias): agnostico al modelo.
\end{itemize}

\begin{lstlisting}[caption={TreeExplainer + KernelExplainer}]
import shap
X_test_df = pd.DataFrame(X_test_sc,
                           columns=features)

# TreeExplainer sobre LightGBM
explainer_tree = shap.TreeExplainer(
    mejor_arbol_obj)
sv_tree = explainer_tree.shap_values(X_test_df)

# Summary bar (importancia global)
shap.summary_plot(sv_tree[1], X_test_df,
    plot_type='bar', feature_names=features)

# Beeswarm (direccion + magnitud)
shap.summary_plot(sv_tree[1], X_test_df,
    feature_names=features)

# Force plot del paciente de mayor riesgo
idx_max = int(np.argmax(
    mejor_arbol_obj.predict_proba(
        X_test_sc)[:,1]))
shap.force_plot(
    explainer_tree.expected_value[1],
    sv_tree[1][idx_max,:],
    X_test_df.iloc[idx_max,:],
    feature_names=features, matplotlib=True)
\end{lstlisting}

\begin{table}[H]
\centering
\caption{Importancia global SHAP -- LightGBM}
\label{tab:shap}
\scriptsize
\begin{tabular}{@{}lrl@{}}
\toprule
\textbf{Variable} & \textbf{SHAP medio} & \textbf{Direccion} \\
\midrule
Glucose       & 0.2847 & Alto$\uparrow$ riesgo \\
BMI           & 0.1423 & Alto$\uparrow$ riesgo \\
Age           & 0.1098 & Mayor$\uparrow$ riesgo \\
DPF           & 0.0921 & Alto$\uparrow$ riesgo \\
Insulin       & 0.0762 & Efecto no lineal \\
Pregnancies   & 0.0614 & Mas$\uparrow$ riesgo \\
SkinThickness & 0.0389 & Efecto pequeno \\
BloodPressure & 0.0287 & Efecto pequeno \\
\bottomrule
\end{tabular}
\end{table}

\textbf{Hallazgos SHAP:} \texttt{Glucose} es el predictor mas influyente;
\texttt{BloodPressure} y \texttt{SkinThickness} tienen impacto marginal,
consistente con el analisis Mann-Whitney U del Paso~5.

%=============================================================
\section{Resultados}
%=============================================================

\subsection{Comparacion global de modelos}

El analisis integral de los diez modelos muestra una progresion clara:

\begin{enumerate}[leftmargin=1.5em, itemsep=0pt]
  \item Modelos base (KNN, NB, SVM, MLP): AUC entre 0.8012 y 0.8243.
  \item Ensambles boosting (XGBoost, LightGBM): AUC entre 0.8354 y 0.8478.
  \item \textbf{Stacking Top-3:} AUC maximo 0.8523.
\end{enumerate}

\subsection{Importancia de variables}

Tanto Mann-Whitney U como SHAP coinciden: \texttt{Glucose} es la
variable mas discriminativa, seguida de \texttt{BMI} y \texttt{Age}.
La consistencia entre ambos metodos valida el hallazgo.

\subsection{Analisis de Lift}

En el primer decil, el Stacking alcanza Lift $\approx 2.92$, significando
que identificar pacientes de mayor riesgo con este modelo es casi tres
veces mas eficiente que hacerlo aleatoriamente.

%=============================================================
\section{Discusion}
%=============================================================

Los resultados son consistentes con la literatura. El ROC AUC de 0.8523
supera los reportados por Sisodia y Sisodia \cite{sisodia2018} (0.80 con
Naive Bayes) y es comparable con los de Rehman et al.\
\cite{rehman2020} (0.84 con Gradient Boosting).

El uso de \textbf{RobustScaler} fue acertado: \texttt{Insulin} y
\texttt{DiabetesPedigreeFunction} presentaban outliers significativos que
habrian sesgado el StandardScaler. La combinacion Winsorizing +
RobustScaler garantizo que KNN, SVM y MLP no fueran afectados por valores
extremos.

La imputacion por mediana de clase aprovecha la informacion del target
sin introducir data leakage. Este enfoque mostro mejoras especialmente
para \texttt{Insulin} (15.2\% de faltantes), donde la diferencia de
medianas entre diabeticos (102 $\mu$U/mL) y no diabeticos (74 $\mu$U/mL)
es sustancial.

Una limitacion del estudio es que el dataset Pima Indians fue recolectado
en una poblacion especifica (mujeres adultas de ascendencia Pima, Arizona,
USA), lo que puede limitar la generalizacion. Trabajos futuros deben
validar el pipeline sobre datasets multiethnicos.

%=============================================================
\section{Conclusiones}
%=============================================================

Este trabajo desarrollo un pipeline completo de Mineria de Datos para
clasificacion de diabetes. Las principales conclusiones son:

\begin{enumerate}[leftmargin=1.5em, itemsep=1pt]
  \item El preprocesamiento exhaustivo (imputacion por mediana de clase,
        Winsorizing, \texttt{log1p} y RobustScaler) es fundamental para el
        desempeno de los modelos, en especial SVM y KNN.
  \item Los modelos de Boosting (LightGBM, XGBoost, Gradient Boosting)
        superan consistentemente a los modelos base, confirmando su
        idoneidad para datos tabulares con distribuciones no normales.
  \item El Stacking con meta-aprendiz y validacion cruzada interna alcanza
        el AUC mas alto (0.8523), demostrando que la combinacion de modelos
        diversos agrega valor predictivo mas alla de los mejores modelos
        individuales.
  \item La interpretabilidad con SHAP revela que \texttt{Glucose},
        \texttt{BMI} y \texttt{Age} son los predictores mas influyentes,
        coherente con la evidencia clinica.
  \item El analisis de Lift confirma la utilidad practica: permite
        priorizar el 10\% de mayor riesgo capturando casi el triple de
        positivos que una seleccion aleatoria.
\end{enumerate}

%=============================================================
\section*{Agradecimientos}
%=============================================================

El autor agradece a la Universidad ESAN y al programa Pronabec Beca 18
por el apoyo academico brindado, y a Ehab Aboelnaga por la publicacion
del dataset en Kaggle.

%=============================================================
\section*{Referencias}
%=============================================================

\begingroup
\renewcommand{\section}[2]{}
\begin{thebibliography}{99}

\bibitem{idf2021}
International Diabetes Federation (IDF),
\textit{IDF Diabetes Atlas}, 10th ed.,
Brussels: IDF, 2021.

\bibitem{who2021}
World Health Organization,
``Diabetes,'' WHO, 2021.
\url{https://www.who.int/news-room/fact-sheets/detail/diabetes}

\bibitem{kaggle2023}
E.\ Aboelnaga, ``Diabetes Dataset,''
Kaggle, 2023.
\url{https://www.kaggle.com/datasets/ehababoelnaga/diabetes-dataset}

\bibitem{kavakiotis2017}
I.\ Kavakiotis et al.,
``Machine Learning and Data Mining Methods in
Diabetes Research,''
\textit{Computational and Structural Biotechnology Journal},
vol.\ 15, pp.\ 104--116, 2017.

\bibitem{sisodia2018}
D.\ Sisodia and D.\ S.\ Sisodia,
``Prediction of Diabetes Using Classification
Algorithms,''
\textit{Procedia Computer Science},
vol.\ 132, pp.\ 1578--1585, 2018.

\bibitem{rehman2020}
A.\ Rehman, S.\ Naz, and I.\ Razzak,
``Leveraging Big Data Analytics in Healthcare,''
\textit{Multimedia Systems},
vol.\ 28, pp.\ 1165--1188, 2022.

\bibitem{chen2016}
T.\ Chen and C.\ Guestrin,
``XGBoost: A Scalable Tree Boosting System,''
in \textit{Proc.\ 22nd ACM SIGKDD},
pp.\ 785--794, 2016.

\bibitem{ke2017}
G.\ Ke et al.,
``LightGBM: A Highly Efficient Gradient
Boosting Decision Tree,''
in \textit{NeurIPS}, vol.\ 30, 2017.

\bibitem{lundberg2017}
S.\ M.\ Lundberg and S.-I.\ Lee,
``A Unified Approach to Interpreting Model
Predictions,''
in \textit{NeurIPS}, vol.\ 30, 2017.

\end{thebibliography}
\endgroup

\end{document}
